\documentclass[12pt,a4paper]{article}

% 基本设置
\usepackage[margin=2.5cm]{geometry} % 页边距
\usepackage{setspace}               % 行距
\usepackage{titlesec}               % 控制标题样式
\usepackage{lipsum}                 % 生成占位文本（可删除）
\usepackage[hidelinks]{hyperref}
\usepackage{enumitem}

\usepackage{luatexja}            % 日语支持
\usepackage{luatexja-fontspec}   % 日文字体与 fontspec 集成

% \usepackage{xeCJK}
% \setCJKmainfont{IPAMincho} % Overleaf 自带日文字体


\setstretch{1.2} % 行距

% 标题格式
\titleformat{\section}{\large\bfseries}{\thesection}{1em}{}
\titleformat{\subsection}{\normalsize\bfseries}{\thesubsection}{1em}{}

\begin{document}

\begin{center}
    \Large \textbf{Mid- and Long-Term Research Plan}
\end{center}

\noindent \textbf{Applicant:} Deyun Lyu \\
\textbf{Affiliation and Position:} Research Fellow, National Institute of Informatics (NII), Tokyo, Japan\\
\textbf{Position Applied For:} Assistant Professor (Informatics and Data Science, Graduate School of Advanced Science and Engineering)\\
\textbf{Field of Specialization:} Intelligent Science (Cyber-Physical Systems, Artificial Intelligence, Software Engineering – with a focus on Quality Assurance)

\section{Overview}

My research and education vision centres on \textbf{advancing systematic methodologies for the \emph{quality assurance} (QA) of intelligent systems} that integrate \emph{artificial intelligence} (AI) with the physical world, while cultivating the next generation of researchers able to address societal challenges. During my doctoral and postdoctoral training, I built a strong foundation in assuring the reliability and trustworthiness of \emph{AI-enabled cyber-physical systems} (AI-CPSs)\hyperlink{paper6}{[6]}\hyperlink{paper7}{[7]}\hyperlink{paper8}{[8]}\hyperlink{paper9}{[9]}, where the AI controller monitors and controls the physical system based on its states and external environments. Building on this, my future research will focus on \textbf{\emph{end-to-end autonomous driving systems}} (E2E-ADS)\hyperlink{paper1}{[1]}\hyperlink{paper2}{[2]}\hyperlink{paper4}{[4]}, \textbf{a leading exemplar} of AI-CPSs. E2E-ADS integrate environmental perception with high-level human driving instructions and directly couple them with control via deep learning models. While this enables advanced driving capabilities, it also introduces unprecedented QA challenges. Their \textbf{black-box nature} and reliance on \textbf{multimodal inputs} (e.g., sensory streams and human driving instructions) make it difficult to ensure robustness, interpretability, and safety under practical conditions. For example, a human instruction such as “turn left at the intersection” may conflict with the visual scene where the left lane is blocked, causing erroneous decisions. Existing QA approaches for conventional AI-CPSs, typically designed for single-modal inputs and modular pipelines, cannot capture the multimodal dependencies of E2E-ADS or diagnose failures in their collapsed black-box architectures. Addressing these challenges requires unified QA methodologies that can holistically assess, explain, and improve system behaviour. In parallel, I will enrich education by supervising student projects, guiding them to publish papers and join competitions, and encouraging involvement in conferences and workshops. My approach emphasizes \textbf{hands-on research} using open-source benchmarks and toolkits, building on my expertise in QA of AI-CPSs.

My ultimate goal is to establish a comprehensive research and education program that extends lessons from E2E-ADS to embodied AI systems (e.g., assistive robots), thereby advancing academic progress, fostering international collaboration, and \textbf{positioning Hiroshima University as a hub for trustworthy embodied AI}.

\section{Mid-Term Research Plan}

Building on the vision outlined above, my mid-term research will address the distinctive QA challenges of E2E-ADS. Unlike conventional AI-CPSs that process structured and unimodal inputs, E2E-ADS integrate high-dimensional sensory streams with human instructions and map them directly to control actions through deep learning. This multimodality introduces cross-modal dependencies and semantic ambiguities that existing QA approaches\hyperlink{paper6}{[6]} cannot capture and solve. To overcome these limitations, I will \textbf{develop systematic QA methodologies} for \textbf{multimodal test generation}, \textbf{test oracle construction}, and \textbf{failure attribution} that explicitly account for multimodal interactions and black-box nature, enabling more robust and interpretable QA of E2E-ADS.

\subsection{Multimodal-Aware Test Generation}

I will \textbf{establish new multimodal-aware testing approaches} tailored to the unique characteristics of E2E-ADS. Conventional testing methods fall short in capturing multimodal dependencies across visual scenes, temporal dynamics, and high-level driving commands. To address this gap, I will design multimodal-aware coverage criteria that can guide the generation of test cases exploring temporal perturbations, cross-modal inconsistencies, and environment–instruction mismatches. These approaches will enable more reliable exploration of system behaviours and help reveal defects that conventional testing approaches might miss.

\subsection{Automated Test Oracle Construction}

Another major challenge in QA of E2E-ADS is the oracle problem: determining the correctness of driving behaviours when multiple valid responses may exist. I will \textbf{construct hierarchical test oracles} that combine metamorphic testing, invariant mining, and specification learning. For instance, metamorphic relations can check the consistency of outputs when sensor data are perturbed but semantics remain unchanged, while invariants mined from repeated runs can capture stable behavioural patterns. Together, these approaches will reduce dependence on expensive manual labelling and make large-scale behavioural evaluation feasible.

\subsection{Failure Attribution and Debugging Support}

When failures occur, developers need actionable insights into whether the problem originates from perception, decision-making, or their cross-modal interactions. Building on my previous work on spectrum-based fault localization (SBFL) for AI-CPSs, I will \textbf{explore methods for isolating minimal failure-inducing subsets of multimodal inputs}, such as specific frames in a visual input sequence or keywords in a human driving instruction. These methods will help diagnose root causes and guide targeted retraining or fine-tuning, bridging the gap between testing and system improvement.

\subsection{Open-Source Frameworks and Benchmarks}

To maximize impact, I will package the proposed methodologies into open-source frameworks validated on benchmarks such as Bench2Drive\hyperlink{paper3}{[3]} and small-scale real-world prototypes. These toolkits will serve as both research artifacts for the scientific community and educational platforms for student projects. In addition, I plan to disseminate the results through workshops and tutorials at international conferences, and to explore collaborations with industry partners to promote technology transfer and real-world adoption.

Through these activities, I aim to publish in top-tier venues in software engineering (e.g., ICSE and TOSEM) and AI (e.g., NeurIPS), while producing reusable tools adopted by academia and industry.

\section{Long-Term Research Plan}

In the long term, I will \textbf{broaden my research from E2E-ADS to embodied AI systems}, such as assistive robots. These systems extend AI-CPSs by combining multimodal perception, natural language interaction, and adaptive action in domains beyond self-driving.

\subsection{Comprehensive Quality Assurance Paradigms}
My research will \textbf{expand from testing to cover the entire lifecycle of assurance}, including runtime monitoring, fault localisation, repair of complex AI components, and safety verification. By closing the loop between offline validation and runtime assurance, I aim to develop frameworks that ensure continuous improvement and dependable operation of embodied AI.

\subsection{Standardization and Industrial Adoption}
I will collaborate with international partners and standardization bodies to \textbf{translate academic advances into widely adopted practices}. This includes contributing to safety guidelines for autonomous driving and embodied AI, as well as establishing open benchmarks that can serve as de facto standards for validation.

\subsection{Interdisciplinary Collaborations}
To address the societal challenges of embodied AI, I will \textbf{pursue interdisciplinary collaborations with researchers} in human-robot interaction, cognitive science, and formal methods. This integration will provide a holistic perspective, ensuring that assurance methodologies are not only technically sound but also aligned with human values and social expectations.

Through these long-term efforts, I aspire to position Hiroshima University as an internationally recognized hub for trustworthy embodied AI research.

\section{Conclusion}

In summary, my mid-term plan will deliver systematic QA methodologies for E2E-ADS, focusing on multimodal test generation, oracle construction, and failure attribution, validated on benchmarks and released as open-source toolkits. In the long term, I will extend these efforts toward comprehensive QA paradigms that integrate testing with runtime monitoring, repair, and safety verification, while contributing to standardization and fostering interdisciplinary collaboration.

On the education side, I will combine hands-on supervision and project-based training in the mid-term with open resources and international programs in the long term. By uniting rigorous research, impactful education, and global collaboration, I aim to \textbf{advance systematic QA of embodied AI systems} and \textbf{establish Hiroshima University as a leading hub for trustworthy intelligent systems}.





 
\section*{References} 
\begin{enumerate}[label={[\arabic*]}, leftmargin=*] 
  \item \hypertarget{paper1}{}Nikkei XTech. ``\href{https://xtech.nikkei.com/atcl/nxt/column/18/03080/012300002/}{自動運転に向けた次世代技術開発が加速}'' (in Japanese). Accessed Sept. 30, 2025.
  \item \hypertarget{paper2}{}B. Jiang et al. ``\href{https://openaccess.thecvf.com/content/ICCV2023/html/Jiang_VAD_Vectorized_Scene_Representation_for_Efficient_Autonomous_Driving_ICCV_2023_paper.html}{VAD: Vectorized Scene Representation for Efficient Autonomous Driving}", ICCV 2023.
  \item \hypertarget{paper3}{}X. Jia et al. ``\href{https://dl.acm.org/doi/10.5555/3737916.3737941}{Bench2Drive: Towards Multi-Ability Benchmarking of Closed-Loop End-to-End Autonomous Driving}", NeurIPS 2025.
  \item \hypertarget{paper4}{}K. Renz et al. ``\href{https://www.computer.org/csdl/proceedings-article/cvpr/2025/436400l993/299arUT3bQA}{SimLingo: Vision-Only Closed-Loop Autonomous Driving with Language-Action Alignment}", CVPR 2025.
  \item \hypertarget{paper5}{}Z. Wang et al. ``\href{https://dl.acm.org/doi/10.1145/3729343}{VLATest: Testing and Evaluating Vision-Language-Action Models for Robotic Manipulation}", FSE 2025.
  \item Z. Zhang, D. Lyu et al. \hypertarget{paper6}{}``\href{https://ieeexplore.ieee.org/document/9844247}{FalsifAI: Falsification of AI-Enabled Hybrid Control Systems Guided by Time-Aware Coverage Criteria}", TSE 2023.
  \item \hypertarget{paper7}{}D. Lyu et al. ``\href{https://dl.acm.org/doi/10.1145/3705307}{SpectAcle: Fault Localisation of AI-Enabled CPS by Exploiting Sequences of DNN Controller Inferences}", TOSEM 2025.
  \item \hypertarget{paper8}{}D. Lyu et al. ``\href{https://www.sciencedirect.com/science/article/abs/pii/S0164121225001438}{Fault Localisation of AI-Enabled Cyber-Physical Systems by Exploiting Temporal Neuron Activation}", JSS 2025.
  \item \hypertarget{paper9}{}D. Lyu et al. ``\href{https://dl.acm.org/doi/10.1145/3638529.3654078}{Search-Based Repair of DNN Controllers of AI-Enabled CPS Guided by System-Level Specifications}", GECCO 2024.
 \end{enumerate}


\end{document}